1992年，总设计师发表南巡讲话，市场化改革的动能重启，长者在抓住时机，统一党内思想，在十四大上干成了来北京的第一件事，确立建设社会主义市场经济，提出“要逐步实行税利分流和分税制”。6月5日，财政部开始在 天津、辽宁、沈阳、大连、浙江、武汉、重庆、青岛和新疆九个试点试行分税制。

1993年，朱镕基正式接手分税制的改革，当时财政收入在国内生产总值的比重从1979年的28.4%降到1993年的12.6%,中央财政在全国财政的比重从46.8%降到31.6%，也就是所谓“双降”的局面。

7月23日，在全国财政会议上，朱镕基坦言，包税制下，中央收入不过40%，支出要占50% , 年年借债，中央财政过不下去了，“不到2000年中央财政就要垮台”。朱镕基撂下狠话：

“这种状况是与市场经济发展背道而驰的，必须调整过来”。

当时由于跟地方约定包税制，所以地方没有交税的积极性，一般都是藏富于民，然后通过收费等其他方式来充实地方金库。

1994年，分税制正式实行，双降的局面一举扭转，到2003年，朱镕基卸任，全国财政占国内GDP的比重从12.6%到18.6%，中央财政从1996年的39%到54.6%。

2011年清华校庆，朱镕基回校参观，借着师生座谈会的机会为自己的政治遗产——分税制辩护。朱镕基的反应非常强烈，指责“攻击分税制，说分税制掏空地方财政，造成农民贫穷的人，“根本就是无知！无知还透顶”。他用2010年的财政数据举例，92年、93年的中央财政比例是28%、27%，而2010年，扣除财政转移支付3万3000亿，中央财政不过是1万5900亿只占20%左右，地方政府有的是钱。

朱镕基认为根子在于当时把买地的钱全都划给了地方政府，而且不纳入预算，所以地方政府开始疯狂卖地。但是他也坦承，分税制也有问题主要在于转移支付，造成了“跑部进钱”。

朱镕基是坦诚的，地方政府财政在民生上的捉襟见肘和土地财政问题 并不是完全由分税制造成的。李 郇、洪国志、黄亮雄考察1999-2008年间，240个地级市地方财政预算内缺口和土地财政的关系发现，地方财政预算内缺口并不能完全解释土地财政的关系。

只有1999-2003年间，随着地方预算内缺口扩大，土地财政的规模才开始在不断扩大；而2003年之后，地方预算内的切口是比较稳定的，但是土地财政的规模却在飞速膨胀。这就说明20003年后，财政压力并不能解释地方对于土地财政的疯狂。

2003年是极为关键的一年，土地财政的起飞更多是因为2002年有两项政策，《招标拍卖挂牌出让国有土地使用权的规定》和 企业和个人所得税改革。前者开了口子，地方政府可以通过土地挂牌拍卖，从市场拿钱；后者就给了地方政府一鞭子，中央和地方以60%：40%的比例来分企业和个人所得税增量。

2001年，地方所得税是1636亿，占地方总收入的21%。被中央拿掉大头之后，地方自然对营业税的依赖进一步上升。 而建筑业又是营业税的第一大户，自此之后，地方政府盯着上房地产，疯狂的发展房地产也就不奇怪了。

其实，省级地方政府角色并没有那么无辜。我国是五级财政制度，由于分税制主要是中央和省级财政之间实行，所以很多的中央的财政转移支付必须通过省级，才能到达市级，再到达市级、县级和乡级，可能会面临层层的盘剥和挪用。

贾康、白景明研究发现，从1994-2000年间，中央的资金集中度从55.7%下降到52.2%，而省级资金集中度却从16.8% 提高到28.8%，年均提高2%。什么意思？就是国家给省里的钱越来越多了，而省里发到县里的钱越来越少了。

2003年，个人所得税改革和土地挂牌拍卖合法化，潘多拉的魔盒正式被打开，土地财政和各个地方政府领导的晋升冲动产生化学反应，固定资产投资飙升，房地产开始起飞，中国经济逐渐开始出现过热的情况。 2003-2006年，每年的 GDP 增长都在10%以上，中央又开始和地方展开新一轮的博弈，宏观调控，为过热的经济降温。

2010年前后，东部开始讲产业升级，腾笼换鸟，中西部就开始蠢蠢欲动。

2007年，郑州就已经成立专门针对富士康的招商工作领导小组，市长亲自任组长。可惜，富士康仍在鼎盛，省级领导时常来亲自拜访，市级领导还常常被拦在门外，郑州并没有任何特殊的待遇，只能时断时续的联系着。2010年，富士康跳楼事件发酵，传出北迁的消息，郑州又兴奋了起来，不过它只是其中的一个，当时还有廊坊、武汉、成都、烟台等许多内陆城市翘首以盼。

最后，郑州给出了极为优厚的条件，五免五减半；降低每年的社保和其他费用1亿美元，享受出口退税。5年之后，当传出中央开始清理税收优惠的消息，郭台铭这回就积极主动地多，还亲自去郑州拜会市长马懿，确认有关50亿的财政补贴的落实。


2012年，三星传出要考虑在海外投资300亿美元建厂， 西安就曾靠着天量的补贴，在北京和重庆的手中，抢下了三星300亿美金的高端储存芯片项目。

在最后2012年4月2日，三星确定项目落户西安的时候，一片哗然，众人纷纷猜测地方政府的补贴力度到底有多强。

媒体一度流传出，2000亿数字的嫁妆：投资30%的补贴，“十免十减半”，修建配套措施，无偿提供土地。 众人惊呼，赔本做买卖。娄勤俭亲自上了凤凰卫视的《神州问答》节目解释，不做亏本买卖，没有“十免十减半”的政策，只是有“五免五减半”，而且是国家的政策，西安只承担运输成本费用。娄省长最后话里还是留了余地，即使不赚钱，

“但是将有260多家配套企业过来，现在就已经有近100多家企业，跟着来投资，这就是一种增值效应，产业带起来了”。

2012年，中央以减税的名义提出营业税改增值税的政策，理由是消除第三产业重复征税，总理亲自赴沪督军。
中央如此劳心费神，自然也别有深意。中国经济的顽疾，是地方政府对于房地产、建筑业的依赖。房地产和建筑业除了能快速拉高 GDP 之外，还是营业税的主要来源，而营业税又是地方政府的主要税种，所以造成了地方政府受到税收激励偏爱房地产。

2016年5月1日，房地产、建筑业等营业税的主体税种纳入营改增，中央和地方的博弈达到高潮。 2016年3-4月，总理在一个月之内，连开了三次会来讲营改增的重要性。

中央拿了营改增拿掉营业税，自然对地方也要所交代，这也就是给地方开的房地产税。房地产税并不是房产税，房产税是一个已经存在的税种，只是一直免征罢了，而房产地产税是指整合房产税和城市土地使用税的新税种。

2017年4月7日，教育部就出台多校划片的政策以朝阳区为试点，规定6月30日之后拿到房产证的，参加多校划区的曾策，希望能够通过增加购买学区房的不确定性来降低房价， 可是刚刚买房的家长，就去教育部门抗议，指责教育部违反义务教育法中，“免试就近入学、学区制和九年一贯对口招生”的规定， 后来多校划片的政策也就不了了之了。

财权与事权相匹配”是分税制的基本原则，有多少钱办多少事。但这种政策有着严格的限制条件，即辖区内经济发展水平的同质化。在中国这样一个大国，各地的经济发展条件和能力并不均衡。2007年，党的十七大报告采纳了“财力与事权相匹配”的提议。

另一句叫做“上面千根线，基层一根针”——上面出政策，下面对口执行，任务最终都压到基层政府，形成“财权层层上收，事权层层下移”的局面。比如“中央点菜，基层埋单”，中央部委下发文件要搞新农合医保、农村危房改造等，地方政府就要掏钱。

“最弱小的一级政府却承担着具有全局意义的支付责任。”天津财经大学教授李炜光评论说。这种体制使县乡两级政府承担了地方的多项公共服务责任。全国两千多个县级组织，曾经有一半以上拖欠教师和离退休人员工资，或大面积拖欠银行贷款和建筑工程款，以至于中央财政亲自出手，出台“三奖一补”政策（指中央对地方缓解县乡财政困难奖励和补助的办法）为此兜底。

“从决策权来看，中国是全世界最集权的国家，从执行权来看，我们又是最分权的国家。”刘尚希说。

在刘尚希看来，造成这一现象的根本原因是，“地方政府只对本级财政负责，它对下没有责任，不愿下移财力，还‘市刮县’向下摊派等”。因此，他开出的改革药方是摈弃“层级财政责任制”，建立“辖区责任制”，即上级政府要对辖区内各级政府的横向和纵向财政平衡负责。比如县一级财政发不出工资，首先应该问责市政府或省政府，而不是由中央财政埋单。
